%============================================================
% ABSTRAK (Bahasa Indonesia)
%   1 paragraf, 150–250 kata, spasi 1, Times New Roman 12.
%   Diakhiri 3–5 kata kunci.
%   Teks di bawah disalin verbatim dari template asli
%   (reference/Template_Skripsi.docx) — ganti dengan abstrak Anda.
%   PENTING: \singlespacing dibungkus \begingroup...\endgroup —
%   tanpa itu deklarasinya bocor dan seluruh isi Bab I–V ikut
%   kehilangan spasi 2.
%============================================================
\begingroup
\singlespacing

\chapter*{ABSTRAK}
\TocHalaman{ABSTRAK}

\noindent
Abstrak merupakan \emph{outline} dari sebuah hasil
penelitian/karya ilmiah/naskah resmi yang memerlukan deskripsi
secara singkat. Abstrak disusun dengan kalimat yang singkat,
jelas, runtut, dan sistematis dan dapat menggambarkan isi laporan
secara keseluruhan. Abstrak disusun dalam bahasa Indonesia,
disusun menjadi 1 paragraf, berkisar antara 150-250 kata, diketik
dengan jarak 1 spasi, menggunakan jenis font Times New Roman,
size 12. Abstrak Skripsi memuat apa dan mengapa penelitian
dikerjakan, bagaimana dikerjakan, dan apa hasil penting yang
diperoleh dari penelitian. Abstrak berisi judul penelitian,
tujuan penelitian, batasan variabel penelitian, serta metode
analisis yang digunakan dalam penelitian. kemudian berisi
ringkasan hasil penelitian dan pembahasan. Selanjutnya,
mengungkapkan kesimpulan dari penelitian. Di bagian bawah
intisari dituliskan kata-kata kunci, bisa berupa kata-kata
penting dalam intisari atau kata yang sering muncul, berjumlah
minimal 3 (tiga) dan maksimal 5 (lima) kata.

\vspace{1em}
\noindent Kata kunci: \KataKunci.

\clearpage
\endgroup
